\section*{ЗАКЛЮЧЕНИЕ}
\label{sec:conclusion}
\addcontentsline{toc}{section}{ЗАКЛЮЧЕНИЕ}

Работата постави въпроса дали една абстракция, конструираща заявките по време на
компилация, може да елиминира цяло семейство от класове инжекции при хетерогенни
хранилища за данни, и къде спира тази гаранция.

\subsection*{Отговор на изследователския въпрос}

Отговорът е положителен, с ясно назовани граници. Когато имената на таблици и
колони и формата на заявката са аргументи на шаблони, а стойностите достигат
заявката само чрез свързани типизирани заместители, недоверен вход не може да
промени структурата на заявката. Това следва от определението на инжекцията като
промяна на синтактичната структура~\cite{su2006essence} и от изискването
аргументите на шаблон да бъдат константни изрази. Аргументът е проверен, а не
само изведен: 15 от 15 опита идентификатор или структура да зависят от стойност
по време на изпълнение са отхвърлени от компилатора по проверяваната причина, при
два компилатора, а при SQLite нито един от 22 приложими полезни товара не е имал
ефект през библиотеката, докато наивната заявка е компрометирана от 4.

Гаранцията обаче не беше пълна в началото на изследването. Стойност, записана
направо в заявката, се приемаше и оставаше несвързана. Това не позволяваше
инжекция, но водеше до тихо грешни резултати. Пропускът е затворен. Именно този
резултат показва най-ясно защо втората част на изследователския въпрос е
равностойна на първата: гаранцията беше изказана като пълна, преди да бъде
измерена.

\subsection*{Приноси}

\begin{enumerate}
    \item \textbf{Научно-приложен принос.} Аргумент, който свързва формалното
          определение на инжекцията с изискването за константни изрази в C++ и
          го разширява от SQL към пет класа инжекции при седем хранилища.
          Установено е, че гаранцията се постига с четири различни механизма и
          че при две от хранилищата липсва текстова граматика, от която полезен
          товар би могъл да излезе.
    \item \textbf{Приложен принос.} Възпроизводим набор за оценка, включен в
          системата за непрекъсната интеграция на библиотеката: 18 компилационни
          теста с контролни случаи и проверка на диагностиката, корпус от 66
          полезни товара с посочен източник за всеки и диференциален тест със
          съзнателно уязвима база за сравнение.
    \item \textbf{Приложен принос.} Открит, измерен и затворен пропуск в
          гаранцията на библиотеката, както и проверка на границата на
          интерфейса, която отхвърля текст на заявка със съобщение, назоваващо
          причината.
    \item \textbf{Отрицателни резултати.} Назовани граници на гаранцията и
          остатъчен риск в неизползван път с по-слаб механизъм, който е
          мъртъв код, но би отслабил гаранцията, ако бъде използван.
\end{enumerate}

\subsection*{Ограничения}

Основното ограничение е, че диференциалният тест е изпълнен само със SQLite. За
останалите шест хранилища гаранцията е потвърдена чрез анализ на изходния код, но
не и чрез измерване. Конекторът за Neo4j не се изпълнява дори от съществуващите
интеграционни тестове на библиотеката. Оракулът вижда само ефекти и не отчита
слепите атаки. Тестът използва една форма на заявка, а корпусът не е изчерпателен.
Най-сетне, самата гаранция има цена: схеми, определяни по време на изпълнение, не
се поддържат, защото именно невъзможността идентификатор да зависи от входа е
същността на защитата.

\subsection*{Бъдеща работа}

Насоките произтичат пряко от ограниченията. Първата е изпълнението на
диференциалния тест при сървърните хранилища. Системата за непрекъсната
интеграция на библиотеката вече стартира PostgreSQL, MySQL, MongoDB, Redis и
Cassandra в контейнери, така че разширяването на теста към тях не изисква нова
инфраструктура. Втората е добавянето на Neo4j към образа за контейнери. Третата
е оракул по двойки, който сравнява две изпълнения и така отчита слепите атаки.
Четвъртата е премахването на мъртвия път с екраниране. Петата е алтернативният
дизайн, при който литералите се свързват вместо да се отхвърлят, оценен по
същия метод. Шестата е самостоятелен анализ на операциите по схемата и на
инжекцията от втори ред, които останаха извън обхвата.
